% BOOK-FILE-ID: f-5cbe4e236eb3
\subsubsection{同じ値を群にまとめる}\label{節：同じ値を群にまとめる}
ガウス記号の値は,中身が少し変わったくらいでは変わらない.
したがって,同じ値を取る項がいくつも続くことになる.
そこで,同じ値を取る項を1つの群にまとめて数えてみよう.
まずは一般項にガウス記号がある数列から見ていく.

\begin{bookpoint}
  \textbf{POINT}\par
  同じ値が続く区間を群とし,各群の項数と値を調べる.
\end{bookpoint}

\noindent\phantomsection\label{証明：ガウス型の群数列}\textbf{【命題1】}\par\nobreak
\begin{smallquote}
  一般項が
  \[
    a_n=\lfloor \sqrt{n} \rfloor\quad(n\ge1)
  \]
  と定義される\bookterm{sequence}{数列}$a_n$は群\bookterm{sequence}{数列}を成す.
\end{smallquote}
\paragraphbf{証明1}
\begin{smallquote}
  正整数$k$について,\,$a_n=k$となる番号$n$は,
  \[
    k\le\sqrt{n}<k+1
  \]
  $n,k\ge1$から,
  \begin{align*}
    & k^2\le n <(k+1)^2                      \\
    & \therefore\quad k^2\le n \le (k+1)^2-1
  \end{align*}
  この範囲にある整数$n$の個数は
  \[
    \{(k+1)^2-1\}-k^2+1=2k+1
  \]
  である.値$k$が$n=k^2$から$2k+1$個並び,次の群では$k+1$となる. \qed
\end{smallquote}

実際の並びも,
\[
  \underbrace{1,1,1}_{3\text{個}},\,
  \underbrace{2,2,2,2,2}_{5\text{個}},\,
  \underbrace{3,3,3,3,3,3,3}_{7\text{個}},\,4,\ldots
\]
となる.

\bookterm{difference}{階差}なら\sectionref{節：階差型}として足す.
\textbf{群の境界と,和の上端が最後の群のどこにあるか}を区別する.

\noindent \bookblocklink{example-gauss-role-2}{problem}{answer}{【例題】}\par\nobreak
\begin{enumerate}[label=(\arabic*),parsep=3pt,beginpenalty=10000,format=\bookitemlink{example-gauss-role-2}{problem}{answer}]
  \item \bookterm{sequence}{数列}を次のように定める.\bookterm{general}{一般項}を求めよ.
    \[
      a_1=2,\qquad
      a_{n+1}=a_n+\lfloor\sqrt{n+2}\rfloor
      \quad(n\ge1).
    \]
    必要ならば,\,$m=\lfloor\sqrt{n+1}\rfloor$を用いてよい.
\end{enumerate}

\noindent \bookblocklink{example-gauss-role-2}{hint}{problem}{【方針】}\par\nobreak
\begin{enumerate}[label=(\arabic*),parsep=3pt,beginpenalty=10000,format=\bookitemlink{example-gauss-role-2}{hint}{problem}]
  \item \bookterm{difference}{階差}を足し,\,$\lfloor\sqrt k\rfloor$の和にします.
    $k=1,2$で値がともに1であることを使い,和の範囲をそろえましょう.
\end{enumerate}

\noindent \bookblocklink{example-gauss-role-2}{answer}{problem}{【解答】}\par\nobreak
\begin{enumerate}[label=(\arabic*),parsep=3pt,beginpenalty=10000,format=\bookitemlink{example-gauss-role-2}{answer}{problem}]
  \item $m=\lfloor\sqrt{n+1}\rfloor$とおくと,
    \[
      a_n=m(n+2)-\frac{m(m+1)(2m+1)}6.
    \]
\end{enumerate}

\noindent \bookblocklink{example-gauss-role-2}{solution}{problem}{【解法】}\par\nobreak
\begin{enumerate}[label=(\arabic*),parsep=3pt,beginpenalty=10000,format=\bookitemlink{example-gauss-role-2}{solution}{problem}]
  \item \bookterm{difference}{階差}を足すと,
    \[
      a_n=2+\sum_{k=3}^{n+1}\lfloor\sqrt k\rfloor
      =\sum_{k=1}^{n+1}\lfloor\sqrt k\rfloor.
    \]\footnote{$n=1$では最初の和を空の和として0とする.}
    $N=n+1,\ m=\lfloor\sqrt N\rfloor$とおく.
    項番号$k$で足し,第$m-1$群までは全部,第$m$群は$N$までを使う.
    \[
      \begin{array}{c|c|c}
        \text{群}&\text{足す項番号}&\text{個数}\\ \hline
        j<m&\substack{j^2\le k\\ k<(j+1)^2}&2j+1\\
        m&\substack{m^2\le k\\ k\le N}&N-m^2+1
      \end{array}
    \]
    最後の群は両端を含むので$N-m^2+1$個.値と個数を掛けて足すと,
    \begin{align*}
      a_n&=\sum_{j=1}^{m-1}j(2j+1)+m(N-m^2+1)\\
      &=\frac{(m-1)m(4m+1)}6\\
      &\quad+m(N-m^2+1)\\
      &=m(N+1)-\frac{m(m+1)(2m+1)}6.
    \end{align*}
    $N=n+1$を戻す.\,$m=1$では最初の和が0となり,初項にも合う.

    \noindent \textbf{【別解】}\\
    $\lfloor\sqrt k\rfloor$は$j^2\le k$を満たす正整数$j$の個数でもある.
    組$(j,k)$の条件$1\le k\le N$, $1\le j$, $j^2\le k$を,\,$j$ごとに数え直す.
    $k=j^2,\ldots,N$の$N-j^2+1$個から,
    \begin{align*}
      a_n&=\sum_{j=1}^{m}(N-j^2+1)\\
      &=m(N+1)-\frac{m(m+1)(2m+1)}6.
    \end{align*}
    群の和と,整数の組の数え直しが同じ答えを与えた.
\end{enumerate}
